\documentclass[10pt]{article}
\usepackage[a4paper,margin=0.65in]{geometry}
\usepackage{amsmath,amssymb,amsfonts,mathtools}
\usepackage[dvipsnames]{xcolor}
\usepackage{array,booktabs,float}
\usepackage[colorlinks=true,linkcolor=blue]{hyperref}
\allowdisplaybreaks
\numberwithin{equation}{section}
\setlength{\jot}{4pt}
\setlength{\emergencystretch}{3em}
\def\mf#1{\mathbb{#1}}
\newcommand{\KK}{\mathcal K}
\newcommand{\Dh}{\hat\Delta}
\newcommand{\OO}{\mathcal O}

\title{Group II, Row IV with $d_\perp$ kept general ($\delta^{mm}=d_\perp$)}
\date{}

\begin{document}
\maketitle

%==========================================================================================
\section{Result}\label{sec:result}
\paragraph{Assumptions.}
\begin{itemize}
\item $d_\perp$ is never set to 2. Every transverse trace is $\delta^{mm}=\delta^{ii}=d_\perp$, both in the polarization sums and in the tensor $z$ integral.
\item The $z$ integral, the only divergent one, is done in $d=2-\epsilon$ dimensions, in $\overline{\rm MS}$, in the scheme used for Group I, Row II and Row V, part 2b.
\item No limit is taken in $\vee$, $k^+$ or $\Lambda$.
\end{itemize}
\paragraph{Notation.} $P\equiv p^++k^+$ and $\xi\equiv k^+/P$, the fraction of the parent carried by the measured gluon. Further,
\begin{equation}
\Delta\equiv y'-y,\qquad \Dh\equiv\frac{\Delta}{|\Delta|},\qquad X\equiv\KK(x-y)=\frac{x-y}{(x-y)^2},\qquad X'\equiv\KK(x'-y'),
\end{equation}
\begin{equation}
\OO\equiv\big[U^{ce'}(y')-U^{ce'}(x')\big]\rho^{e'}(x')\,\big[U^{ce}(y)-U^{ce}(x)\big]\rho^{e}(x)
\end{equation}
(with $U^{bc}=U^{cb}$ this is your $[U^{e'c}(y')-U^{e'c}(x')][U^{ec}(y)-U^{ce}(x)]\rho^{e'}(x')\rho^e(x)$), and
\begin{equation}
E\equiv\frac{2}{\epsilon}+2\gamma_E+\log\frac{\Delta^2\mu^2_{\overline{MS}}}{4}=-\Big[\frac{1}{\bar\epsilon}-\log\frac{\Delta^2\mu^2_{\overline{MS}}e^{2\gamma_E}}{4}\Big]\equiv-\bar E .
\label{eq:E}
\end{equation}
The pole is infrared ($\epsilon<0$); $\frac1{\bar\epsilon}=-\frac2\epsilon$ is your convention.
\paragraph{Result.} The splitting function is
\begin{equation}
\hat\Phi(\xi)\equiv\xi(1-\xi)+\frac{\xi}{1-\xi}+\frac{1-\xi}{\xi},
\end{equation}
and Row IV is
\begin{equation}
\boxed{\begin{aligned}
\frac{d^3N_{\rm IV}}{d^3k}={}&\frac{1}{(2\pi)^3}\frac{g^4N_c}{16\pi^4}\frac{1}{k^+}\int_{k^+/\vee}^{k^+/(k^++\Lambda)}\frac{d\xi}{\xi^2}\int_{x,x',y,y'}e^{-ik\cdot\Delta/\xi}\\
&\times\Big\{X\cdot X'\Big[d_\perp\,\hat\Phi(\xi)\,\bar E-2\Big(\frac{\xi}{1-\xi}+\frac{1-\xi}{\xi}\Big)\Big]-2\,d_\perp\,\xi(1-\xi)\,(X\cdot\Dh)(X'\cdot\Dh)\Big\}\,\OO .
\end{aligned}}
\label{eq:result}
\end{equation}
\begin{itemize}
\item $d_\perp$ enters in two places: it multiplies the whole pole term $\hat\Phi\bar E$, and the $(X\cdot\Dh)(X'\cdot\Dh)$ term.
\item The term $-2(\frac{\xi}{1-\xi}+\frac{1-\xi}{\xi})$ does not depend on $d_\perp$. It comes from the traceless part of the $z$ integral, contracted with $\delta_{mm'}$ (Sec.~\ref{sec:contract}).
\item The term $\frac{2}{k^+}(\delta_{mk'}\delta_{m'k}-\delta_{m'k'}\delta_{km})$ of your bracket does not contribute, for any $d_\perp$.
\end{itemize}
In $E$ form, Eq.~\eqref{eq:result} is $-\frac{1}{(2\pi)^3}\frac{g^4N_c}{16\pi^4}\frac1{k^+}\int\frac{d\xi}{\xi^2}\int e^{-ik\cdot\Delta/\xi}\{X\cdot X'[d_\perp\hat\Phi E+2(\frac{\xi}{1-\xi}+\frac{1-\xi}{\xi})]+2d_\perp\xi(1-\xi)(X\cdot\Dh)(X'\cdot\Dh)\}\OO$.

\paragraph{Soft endpoint ($\xi\to1$).} Write $\xi_0=\frac{k^+}{\vee}$, $\xi_1=\frac{k^+}{k^++\Lambda}$ and $g(\xi)=\xi^{-2}e^{-ik\cdot\Delta/\xi}$. Exactly,
\begin{align}
\frac{d^3N_{\rm IV}}{d^3k}={}&\frac{1}{(2\pi)^3}\frac{g^4N_c}{16\pi^4}\frac{1}{k^+}\Big[\log\frac{\vee}{\Lambda}-\log\frac{\vee}{k^+}+\log\Big(1-\frac{k^+}{\vee}\Big)+\log\Big(1+\frac{\Lambda}{k^+}\Big)\Big]\int_{x,x',y,y'}e^{-ik\cdot\Delta}\,X\cdot X'\,\big(d_\perp\bar E-2\big)\,\OO\notag\\
&+\frac{1}{(2\pi)^3}\frac{g^4N_c}{16\pi^4}\frac{1}{k^+}\int_{\xi_0}^{\xi_1}d\xi\int_{x,x',y,y'}\Big\{X\cdot X'\Big[\big(d_\perp\bar E-2\big)\frac{g(\xi)-g(1)}{1-\xi}\notag\\
&\qquad+g(\xi)\Big(d_\perp\bar E\Big(\xi(1-\xi)-1+\frac{1-\xi}{\xi}\Big)-2\Big(\frac{1-\xi}{\xi}-1\Big)\Big)\Big]-2d_\perp\xi(1-\xi)\,g(\xi)\,(X\cdot\Dh)(X'\cdot\Dh)\Big\}\OO .
\label{eq:endpoint}
\end{align}
The coefficient of the endpoint logarithm is $d_\perp\bar E-2$.

\paragraph{Pole and LO.} The LO cross section (emitter at $y$) is
\begin{equation}
\frac{d^3N_{\rm LO}}{d^3k}(k^+,\mathbf k)=\frac{1}{(2\pi)^3}\frac{g^2}{2\pi^2k^+}\int_{x,x',y,y'}e^{-ik\cdot\Delta}X\cdot X'\,\OO .
\end{equation}
With it, the pole of Eq.~\eqref{eq:result} is
\begin{equation}
\frac{d^3N_{\rm IV}}{d^3k}\bigg|_{\rm pole}=\frac{\alpha_s}{2\pi}\,\frac1{\bar\epsilon}\int_{\xi_0}^{\xi_1}\frac{d\xi}{\xi^3}\,\frac{d_\perp}{2}\,P^{(r)}_{gg}(\xi)\,\frac{d^3N_{\rm LO}}{d^3k}\Big(\frac{k^+}{\xi},\frac{\mathbf k}{\xi}\Big),\qquad P^{(r)}_{gg}(\xi)=2N_c\hat\Phi(\xi).
\label{eq:pole}
\end{equation}
So the residue is $\frac{d_\perp}{2}P^{(r)}_{gg}$: the polarization trace multiplies the real part of the splitting function.

\paragraph{The two usual choices.}
\begin{itemize}
\item $d_\perp=2-\epsilon$ (conventional dimensional regularization, as in Group I Rows II and V). Then $d_\perp\bar E=2\bar E+2$, and the braces of \eqref{eq:result} become $2\{X\cdot X'[\hat\Phi\bar E+\xi(1-\xi)]-2\xi(1-\xi)(X\cdot\Dh)(X'\cdot\Dh)\}$. This is the result I gave before, in the previous version of \texttt{Group\_II\_corrected.pdf}, where $d_\perp=2-\epsilon$ had been inserted.
\begin{itemize}
\item Up to the corrections given earlier (the constant in the logarithm, the $\xi$ limits), it also agrees with your Eq.~(1.15): there the ``$+1$'' in $[\frac1{\bar\epsilon}+1-\log]$ is $d_\perp\bar E$ with $d_\perp=2-\epsilon$, while the last line keeps $d_\perp$. Note that $\int_1^{k^+/\vee}=-\int_{k^+/\vee}^{1}$.
\item The endpoint coefficient becomes $2\bar E$.
\end{itemize}
\item $d_\perp=2$ (polarizations in two dimensions). The braces become $2\{X\cdot X'[\hat\Phi\bar E-\frac{\xi}{1-\xi}-\frac{1-\xi}{\xi}]-2\xi(1-\xi)(X\cdot\Dh)(X'\cdot\Dh)\}$ and the endpoint coefficient is $2\bar E-2$.
\end{itemize}
The difference between the two choices is finite, $2\hat\Phi\,X\cdot X'$ in the braces. The same choice of $d_\perp$ must be used in every row, so that the poles and the finite parts match between rows and with the counterterm.

%==========================================================================================
\section{Derivation}\label{sec:derivation}
\subsection{Step 1: the row}
Row IV is the term $4\,\mathfrak C'^\dagger\mathfrak C'$ of Eq.~(CS2). The parent gluon $P$ is emitted by the charge at $x$ ($x'$ in the conjugate amplitude). It crosses the target at $y$ ($y'$) and splits after the target into the measured gluon $k^+$ at $w$ ($w'$) and the unobserved gluon $p^+$ at $z$. Inserting
\begin{itemize}
\item $\mf A^{(1)}(P,-y)=\frac{ig}{\sqrt2\pi\sqrt P}\int_x\KK(x-y)\rho(x)$,
\item and $\mf C_1$, with vertex bracket
\begin{equation}
V_{ij;km}\equiv\delta_{km}\delta_{ij}-\frac{P}{p^+}\delta_{jm}\delta_{ik}-\frac{P}{k^+}\delta_{im}\delta_{kj}
\end{equation}
($i$: measured gluon, $j$: unobserved gluon, $k$: parent, $m$: index of the splitting kernel),
\end{itemize}
gives your Eq.~(1.12):
\begin{align}
\frac{d^3N}{d^3k}={}&\frac{1}{(2\pi)^3}\frac{g^4}{4\pi^4}\int_\Lambda^{\vee-k^+}\frac{dp^+}{2\pi}\int_{z,w,w'}e^{-ik\cdot(w'-w)}\int_{x,x',y,y'}\KK^m(w-z)\KK^k(x-y)\KK^{m'}(w'-z)\KK^{k'}(x'-y')\frac{p^+k^+}{P^4}\notag\\
&\times\delta^{(d)}\Big[y-w+\frac{p^+}{P}(w-z)\Big]\delta^{(d)}\Big[y'-w'+\frac{p^+}{P}(w'-z)\Big]\,V_{ij;km}V_{ij;k'm'}\times[\text{colour}].
\label{eq:start}
\end{align}
The prefactor is $4\times|\mf A^{(1)}(P)|^2\times|\mf C_1|^2$: $4\cdot\frac{g^2}{2\pi^2P}\cdot\frac{g^2}{8\pi^2P}\frac{p^+k^+}{P^2}=\frac{g^4}{4\pi^4}\frac{p^+k^+}{P^4}$.

\subsection{Step 2: the polarization sum (\texorpdfstring{$\delta^{ii}=d_\perp$}{delta ii = d}).}
The measured and the unobserved gluon are summed over $i$ and $j$. The nine products give:
\begin{center}\small
\begin{tabular}{@{}lll@{}}\toprule
product & $\sum_{ij}$ & result\\\midrule
$\delta_{km}\delta_{ij}\cdot\delta_{k'm'}\delta_{ij}$ & $\delta_{ij}\delta_{ij}=d_\perp$ & $d_\perp\,\delta_{km}\delta_{k'm'}$\\
$\delta_{km}\delta_{ij}\cdot(-\frac{P}{p^+})\delta_{jm'}\delta_{ik'}$ and its mirror & & $-2\frac{P}{p^+}\delta_{km}\delta_{k'm'}$\\
$\delta_{km}\delta_{ij}\cdot(-\frac{P}{k^+})\delta_{im'}\delta_{k'j}$ and its mirror & & $-2\frac{P}{k^+}\delta_{km}\delta_{k'm'}$\\
$(\frac{P}{p^+})^2\delta_{jm}\delta_{ik}\delta_{jm'}\delta_{ik'}$ & & $(\frac{P}{p^+})^2\delta_{mm'}\delta_{kk'}$\\
$(\frac{P}{k^+})^2\delta_{im}\delta_{kj}\delta_{im'}\delta_{k'j}$ & & $(\frac{P}{k^+})^2\delta_{mm'}\delta_{kk'}$\\
$\frac{P^2}{p^+k^+}\delta_{jm}\delta_{ik}\delta_{im'}\delta_{k'j}$ and its mirror & & $2\frac{P^2}{p^+k^+}\delta_{mk'}\delta_{km'}$\\\bottomrule
\end{tabular}
\end{center}
Only the first product contains a trace. Multiplying by $\frac{p^+k^+}{P^4}$, with $p^++k^+=P$:
\begin{equation}
\frac{p^+k^+}{P^4}\sum_{ij}V_{ij;km}V_{ij;k'm'}=\Big[\frac{d_\perp p^+k^+}{P^4}-\frac{2}{P^2}\Big]\delta_{km}\delta_{k'm'}+\Big[\frac{k^+}{p^+P^2}+\frac{p^+}{k^+P^2}\Big]\delta_{mm'}\delta_{kk'}+\frac{2}{P^2}\,\delta_{mk'}\delta_{km'}.
\label{eq:polsum}
\end{equation}
This was checked numerically for $d_\perp=2,3,4,5$ (@@POL@@).

\subsection{Step 3: the \texorpdfstring{$w$}{w} and \texorpdfstring{$w'$}{w'} integrations in \texorpdfstring{$d_\perp$}{d} dimensions}
In $d$ transverse dimensions:
\begin{itemize}
\item the splitting coefficient $\mf C_1$ carries $\delta^{(d)}$;
\item its kernel is the $d$-dimensional Fourier transform of $\mathbf q^m/\mathbf q^2$:
\begin{equation}
\KK^m_d(r)\equiv-2\pi i\int\frac{d^dq}{(2\pi)^d}e^{iq\cdot r}\frac{q^m}{q^2}=\pi^{1-\frac d2}\,\Gamma\Big(\frac d2\Big)\frac{r^m}{|r|^d},\qquad \KK_2(r)=\KK(r).
\label{eq:Kd}
\end{equation}
\end{itemize}
The argument of the first $\delta$ is $\frac{k^+}{P}\big(w_0-w\big)$ with $w_0=\frac{P}{k^+}y-\frac{p^+}{k^+}z$. Hence
\begin{itemize}
\item $\int d^dw\,\delta^{(d)}[\dots]=(\frac{P}{k^+})^d$;
\item $w-z=\frac{P}{k^+}(y-z)$, and $\KK_d$ has degree $1-d$, so $\KK_d(w-z)=(\frac{P}{k^+})^{1-d}\KK_d(y-z)$.
\end{itemize}
Per side the factor is $(\frac{P}{k^+})^d(\frac{P}{k^+})^{1-d}=\frac{P}{k^+}$, so both sides give $(\frac{P}{k^+})^2$ \emph{for any $d$}. In particular there is no $(\frac{P}{k^+})^{-\epsilon}$, which would change the finite part. The phase is $e^{-ik\cdot(w'-w)}=e^{-ik\cdot\Delta\,P/k^+}=e^{-ik\cdot\Delta/\xi}$.

Eq.~\eqref{eq:polsum} times $(\frac{P}{k^+})^2$, with $\frac1{k^+}$ taken out, is your bracket of Eq.~(1.13):
\begin{equation}
\frac{1}{k^+}\Big\{\underbrace{\Big[\frac{d_\perp p^+}{P^2}-\frac{2}{k^+}\Big]}_{c_1}\delta_{km}\delta_{k'm'}+\underbrace{\Big[\frac{1}{p^+}+\frac{p^+}{k^{+2}}\Big]}_{c_2}\delta_{mm'}\delta_{kk'}+\underbrace{\frac{2}{k^+}}_{c_3}\delta_{mk'}\delta_{km'}\Big\}.
\label{eq:bracket}
\end{equation}

\subsection{Step 4: the Wilson lines}
As in your Eq.~(1.14), and checked before (\texttt{Groups\_II\_III\_check.pdf}), the colour factor reduces to $N_c\OO$. So
\begin{equation}
\frac{d^3N}{d^3k}=\frac{1}{(2\pi)^3}\frac{g^4N_c}{4\pi^4}\frac{1}{k^+}\int_\Lambda^{\vee-k^+}\frac{dp^+}{2\pi}\int_{x,x',y,y'}e^{-ik\cdot\Delta/\xi}\,X^kX'^{k'}\big[c_1\delta_{km}\delta_{k'm'}+c_2\delta_{mm'}\delta_{kk'}+c_3\delta_{mk'}\delta_{km'}\big]J^{mm'}\,\OO ,
\label{eq:beforez}
\end{equation}
with the $z$ integral
\begin{equation}
J^{mm'}\equiv\int d^dz\,\KK^m_d(y-z)\,\KK^{m'}_d(y'-z)\quad\times\Big(\frac{\mu^2e^{\gamma_E}}{4\pi}\Big)^{\epsilon/2}.
\end{equation}

\subsection{Step 5: the tensor \texorpdfstring{$z$}{z} integral in \texorpdfstring{$d$}{d} dimensions}\label{sec:J}
With Eq.~\eqref{eq:Kd}, $\int d^dz\,e^{-i(q+q')\cdot z}=(2\pi)^d\delta^{(d)}(q+q')$ gives
\begin{equation}
J^{mm'}=(-2\pi i)^2\int\frac{d^dq}{(2\pi)^d}\frac{q^m(-q^{m'})}{q^4}\,e^{iq\cdot(y-y')}=4\pi^2\int\frac{d^dq}{(2\pi)^d}\frac{q^mq^{m'}}{q^4}\,e^{-iq\cdot\Delta}.
\end{equation}
Use
\begin{itemize}
\item $\int\frac{d^dq}{(2\pi)^d}\frac{e^{iq\cdot r}}{(q^2)^2}=\frac{\Gamma(\frac d2-2)}{16\pi^{d/2}}(r^2)^{2-\frac d2}$;
\item $q^mq^{m'}\to-\partial^m\partial^{m'}$;
\item $\partial^m\partial^{m'}(r^2)^a=2a\,\delta^{mm'}(r^2)^{a-1}+4a(a-1)\,r^mr^{m'}(r^2)^{a-2}$;
\item and, with $a=2-\frac d2$, $a\Gamma(\frac d2-2)=-\Gamma(\frac d2-1)$ and $a(a-1)\Gamma(\frac d2-2)=\Gamma(\frac d2)$.
\end{itemize}
Then
\begin{equation}
J^{mm'}=\frac{\pi^{2-\frac d2}}{2}\,(\Delta^2)^{1-\frac d2}\Big(\frac{\mu^2e^{\gamma_E}}{4\pi}\Big)^{\epsilon/2}\Big[\Gamma\Big(\frac d2-1\Big)\,\delta^{mm'}-2\,\Gamma\Big(\frac d2\Big)\,\Dh^m\Dh^{m'}\Big]\equiv A\,\delta^{mm'}+B\,\Dh^m\Dh^{m'}.
\label{eq:J}
\end{equation}
The integral converges for $2<d<4$ ($\epsilon<0$): the pole is infrared. With $d=2-\epsilon$:
\begin{align}
A&=\frac{\pi}{2}\Big(\frac{\Delta^2\mu^2e^{\gamma_E}}{4}\Big)^{\epsilon/2}\Gamma\Big(-\frac\epsilon2\Big)=-\frac{\pi}{2}E+O(\epsilon),\\
B&=-\pi\Big(\frac{\Delta^2\mu^2e^{\gamma_E}}{4}\Big)^{\epsilon/2}\Gamma\Big(1-\frac\epsilon2\Big)=-\pi+O(\epsilon).
\end{align}
The trace is $J^{mm}=d_\perp A+B$; with $d_\perp=d$ it equals $\pi^{2-\frac d2}\Gamma(\frac d2-1)(\Delta^2)^{1-\frac d2}(\dots)^{\epsilon/2}$, the scalar integral of Group I (Row II, Row V part 2b).

Numerical checks of Eq.~\eqref{eq:J} by quadrature:
\begin{itemize}
\item $d=2$: the traceless part $J^{11}-J^{22}$ for $\Delta\parallel\hat e_1$ is @@J2@@ (closed form $B=-\pi$);
\item $d=3$: the trace is @@J3T@@ (closed form $\pi/|\Delta|=$ @@J3TC@@) and $\Dh\cdot J\cdot\Dh=$ @@J33@@ (closed form $0$).
\end{itemize}

\subsection{Step 6: the contraction (\texorpdfstring{$\delta^{mm}=d_\perp$}{delta mm = d})}\label{sec:contract}
\begin{itemize}
\item $c_1\,\delta_{km}\delta_{k'm'}$: $X^kX'^{k'}J^{kk'}=A\,X\cdot X'+B\,(X\cdot\Dh)(X'\cdot\Dh)$.
\item $c_3\,\delta_{mk'}\delta_{km'}$: $X^kX'^{k'}J^{k'k}$. This is the same, because $J$ is symmetric.
\item $c_2\,\delta_{mm'}\delta_{kk'}$: $X\cdot X'\,J^{mm}=X\cdot X'\,(d_\perp A+B)$.
\end{itemize}
Since $c_1+c_3=\frac{d_\perp p^+}{P^2}$, the $\frac{2}{k^+}$ terms of your bracket cancel for any $d_\perp$:
\begin{equation}
\Sigma\equiv\frac{d_\perp p^+}{P^2}\Big[A\,X\cdot X'+B\,(X\cdot\Dh)(X'\cdot\Dh)\Big]+\Big[\frac1{p^+}+\frac{p^+}{k^{+2}}\Big]\big(d_\perp A+B\big)\,X\cdot X'.
\label{eq:Sigma}
\end{equation}
This is exact in $d$. Inserting $A$ and $B$:
\begin{equation}
\Sigma=-\frac\pi2\Big\{X\cdot X'\Big[d_\perp\Big(\frac{p^+}{P^2}+\frac1{p^+}+\frac{p^+}{k^{+2}}\Big)E+2\Big(\frac1{p^+}+\frac{p^+}{k^{+2}}\Big)\Big]+2\,d_\perp\frac{p^+}{P^2}(X\cdot\Dh)(X'\cdot\Dh)\Big\}+O(\epsilon).
\end{equation}
The ``$+2$'' is $B$ in the trace $J^{mm}$. The $O(\epsilon)$ terms of $A$ and $B$ multiply $d_\perp$, which is finite, so they drop.

\subsection{Step 7: \texorpdfstring{$p^+\to\xi$}{p+ to xi}}
We have $p^+=k^+\frac{1-\xi}{\xi}$, $dp^+=-k^+\frac{d\xi}{\xi^2}$, and the limits $\Lambda<p^+<\vee-k^+$ become $\frac{k^+}{\vee}<\xi<\frac{k^+}{k^++\Lambda}$. Then
\begin{equation}
k^+\Big(\frac{p^+}{P^2}+\frac1{p^+}+\frac{p^+}{k^{+2}}\Big)=\hat\Phi(\xi),\qquad k^+\Big(\frac{1}{p^+}+\frac{p^+}{k^{+2}}\Big)=\frac{\xi}{1-\xi}+\frac{1-\xi}{\xi},\qquad k^+\frac{p^+}{P^2}=\xi(1-\xi).
\end{equation}
Collecting the prefactors, $\frac{g^4N_c}{4\pi^4}\cdot\frac{1}{2\pi}\cdot\big(-\frac{\pi}{2}\big)=-\frac{g^4N_c}{16\pi^4}$. With $E=-\bar E$ this gives Eq.~\eqref{eq:result}.

\subsection{Step 8: the endpoint}
\begin{itemize}
\item The only singular coefficient at $\xi\to1$ is $\frac{\xi}{1-\xi}=\frac{1}{1-\xi}-1$, in $\hat\Phi$ and in the $d_\perp$-independent term. Together it multiplies $(d_\perp E+2)\,X\cdot X'$.
\item Exactly, $\int_{\xi_0}^{\xi_1}d\xi\frac{h(\xi)}{1-\xi}=\int_{\xi_0}^{\xi_1}d\xi\frac{h(\xi)-h(1)}{1-\xi}+h(1)\log\frac{1-\xi_0}{1-\xi_1}$.
\item The limits give
\begin{equation}
\log\frac{1-\xi_0}{1-\xi_1}=\log\frac{(\vee-k^+)(k^++\Lambda)}{\vee\Lambda}=\log\frac{\vee}{\Lambda}-\log\frac{\vee}{k^+}+\log\Big(1-\frac{k^+}{\vee}\Big)+\log\Big(1+\frac{\Lambda}{k^+}\Big).
\end{equation}
\end{itemize}
This gives Eq.~\eqref{eq:endpoint}. It was checked numerically for arbitrary $d_\perp$ and $\bar E$ (@@END@@).

\subsection{Step 9: the pole as a convolution with LO}
\begin{itemize}
\item $\frac{d^3N_{\rm LO}}{d^3k}(\frac{k^+}{\xi},\frac{\mathbf k}{\xi})$ has the prefactor $\frac{g^2\xi}{2\pi^2k^+}$ and the phase $e^{-ik\cdot\Delta/\xi}$.
\item So $\int e^{-ik\cdot\Delta/\xi}X\cdot X'\,\OO=(2\pi)^3\frac{2\pi^2k^+}{g^2\xi}\frac{d^3N_{\rm LO}}{d^3k}(\frac{k^+}{\xi},\frac{\mathbf k}{\xi})$.
\item The $\frac1{\bar\epsilon}$ part of Eq.~\eqref{eq:result} is then $\frac{g^2N_c}{8\pi^2}\frac{1}{\bar\epsilon}\int\frac{d\xi}{\xi^3}d_\perp\hat\Phi\,\frac{d^3N_{\rm LO}}{d^3k}$, which is Eq.~\eqref{eq:pole} with $\frac{g^2N_c}{8\pi^2}=\frac{\alpha_sN_c}{2\pi}$.
\item The measure is $d\xi/\xi^3$.
\end{itemize}

%==========================================================================================
\section{Checks}
Script \texttt{check\_d.py}, results in \texttt{check\_d.json}.
\begin{center}\small
\begin{tabular}{@{}p{0.55\textwidth}p{0.38\textwidth}@{}}\toprule
check & result\\\midrule
polarization sum \eqref{eq:polsum}, $d_\perp=2,\dots,5$ & @@POL@@\\
$z$ integral \eqref{eq:J}, $d=2$ traceless part & @@J2@@ vs.\ $-\pi$\\
$z$ integral \eqref{eq:J}, $d=3$ trace and $\Dh J\Dh$ & rel.\ @@J3R@@\\
contraction \eqref{eq:Sigma} with $J=A\delta+B\Dh\Dh$, $J^{mm}=d_\perp A+B$ (\texttt{sympy}) & @@SYM@@\\
$d_\perp=2-\epsilon$: reduces to the earlier result (\texttt{sympy}) & @@RED@@\\
$\xi$ forms of the coefficients (\texttt{sympy}) & @@XI@@\\
endpoint decomposition \eqref{eq:endpoint}, exact limits & rel.\ @@END@@\\\bottomrule
\end{tabular}
\end{center}
\end{document}
